% Fichier généré automatiquement par tools/generate_corrections.py.
% Source : DYN/DYN-03-Inertie/40_Parallelepipede/40_Parallelepipede.tex
% Statut : complet (2/2 question(s) renseignée(s), 0 reprise(s)).
\begin{corrigebox}[Corrigé extrait de la banque]
\CorrigeQuestion{1}
Pour des raisons de symétrie, on a directement $\vect{OG}=\dfrac{a}{2}\vect{x}+\dfrac{b}{2}\vect{y}+\dfrac{c}{2}\vect{z}$.
\CorrigeQuestion{2}
Notons (1) le parallélépipède rectangle et (2) le cylindre (plein). On note $\mathcal{B}_0=\base{x}{y}{z}$
On a $\inertie{G}{1}=\matinertie{A_1}{B_1}{C_1}{0}{0}{0}{\mathcal{B}_0}$ et
$\inertie{G}{2}=\matinertie{A_2}{B_2}{A_2}{0}{0}{0}{\mathcal{B}_0}$ (attention l'axe du cylindre est $\vect{y}$).

On a donc $\inertie{G}{S}=\matinertie{A_1-A_2}{B_1-B_2}{C_1-A_2}{0}{0}{0}{\mathcal{B}_0}$.

Par ailleurs, $m=m_1-m_2$
et $\vect{AG}=\dfrac{b}{2}\vect{y}$; donc $\inertie{A}{S}=\matinertie{A_1-A_2+m\dfrac{b^2}{4}}{B_1-B_2}{C_1-A_2+m\dfrac{b^2}{4}}{0}{0}{0}{\mathcal{B}_0}$.

Enfin, $\vect{OG}=\dfrac{a}{2}\vect{x}+\dfrac{b}{2}\vect{y}+\dfrac{c}{2}\vect{z}$; donc
 $\inertie{O}{S}=
\matinertie{A_1-A_2}{B_1-B_2}{C_1-A_2}{0}{0}{0}{\mathcal{B}_0}
+m\matinertie{\dfrac{b^2}{4}+\dfrac{c^2}{4}}{\dfrac{a^2}{4}+\dfrac{c^2}{4}}{\dfrac{a^2}{4}+\dfrac{b^2}{4}}{-\dfrac{bc}{4}}{-\dfrac{ac}{4}}{-\dfrac{ab}{4}}{\mathcal{B}_0}$.
\end{corrigebox}
